% ══════════════════════════════════════════════════════════════════
% Source du CV téléchargeable depuis l'accueil (français, 1 page, police Charter)
% Version française généraliste (novembre 2026), même mise en page que le CV Salomon.
% Compiler : pdflatex CV_Gauthier_Bonhomme_FR.tex  (ou tectonic)
% Le PDF produit remplace Gauthier_Bonhomme_CV.pdf à la racine du site.
% L'ancien CV anglais (mécatronique / IA) reste dans CV_Gauthier_Bonhomme_EN.tex / .pdf.
% ══════════════════════════════════════════════════════════════════
\documentclass[10pt,a4paper]{article}
\usepackage[top=1.0cm,bottom=0.9cm,left=1.3cm,right=1.3cm]{geometry}
\usepackage[T1]{fontenc}
\usepackage[utf8]{inputenc}
\usepackage[french]{babel}
\frenchsetup{StandardLists=true}
\usepackage{charter}
\usepackage{titlesec}
\usepackage{needspace}
\usepackage{enumitem}
\usepackage{xcolor}
\usepackage{qrcode}
\definecolor{linkblue}{RGB}{21,76,155}
\usepackage[colorlinks=true,urlcolor=linkblue,linkcolor=linkblue,pdftitle={Gauthier Bonhomme - CV},pdfauthor={Gauthier Bonhomme}]{hyperref}

\pagestyle{empty}
\setlength{\parindent}{0pt}
\setlength{\parskip}{0.1em}

% titres de section : bleu, gras, filet bleu
\titleformat{\section}{\needspace{4\baselineskip}\large\bfseries\color{linkblue}}{}{0pt}{}[\vspace{-0.75em}{\color{linkblue}\rule{\textwidth}{0.5pt}}]
\titlespacing{\section}{0pt}{1.0em}{0.55em}

\newcommand{\lk}{\,\textsuperscript{$\nearrow$}}
\newcommand{\link}[2]{\href{#1}{#2\lk}}
\newcommand{\sep}{\,{\scriptsize\textbullet}\,}
% entrée : titre à gauche, date à droite
\newcommand{\entry}[2]{\par\vspace{0.5em}#1\hfill #2\par}

\setlist[itemize]{leftmargin=1.2em,itemsep=0.05em,parsep=0pt,topsep=0.1em,label={\color{linkblue}\scriptsize\textbullet}}

\begin{document}

% ───── en-tête (QR code vers le portfolio à droite) ─────
\begin{minipage}[c]{0.13\linewidth}\hfill\end{minipage}%
\begin{minipage}[c]{0.74\linewidth}
\centering
  {\Huge\bfseries Gauthier Bonhomme}\\[0.45em]
  Élève ingénieur en dernière année --- \textbf{Arts et Métiers (ENSAM)} --- Aix-en-Provence\\[0.15em]
  {\color{linkblue}\bfseries Stage de fin d'études}\sep 6 mois dès début 2027\\[0.25em]
  {\small +33 6 08 27 62 16\sep \href{mailto:gauthier.bonhomme30@gmail.com}{gauthier.bonhomme30@gmail.com}\sep permis B et A}\\[0.1em]
  {\small\color{linkblue}$\rightarrow$ Rapports, photos et albums :\ \textbf{\href{https://gauthierbonhomme.com}{gauthierbonhomme.com}}}
\end{minipage}%
\begin{minipage}[c]{0.13\linewidth}
\raggedleft
{\color{linkblue}\qrcode[height=1.55cm]{https://gauthierbonhomme.com}}\\[0.15em]
{\scriptsize\color{linkblue}\bfseries Portfolio\hspace{0.32cm}}
\end{minipage}

% ───── expérience ─────
\section{Expérience}

\entry{\textbf{Royal Canin} --- stagiaire, usine d'Aimargues (Gard)}{juin -- août 2026}
\begin{itemize}
  \item Suivi du chantier de construction de l'infirmerie du site, sur le terrain avec les équipes.
  \item Amélioration du logiciel de maintenance de l'usine, au quotidien aux côtés des équipes de maintenance.
\end{itemize}

\entry{\textbf{Hall b} --- équipe fondatrice d'une salle de sport, Saint-Dionisy (Gard)}{juin -- août 2025}
\begin{itemize}
  \item Lancement de A à Z : coordination des artisans et suivi du planning des travaux d'aménagement.
  \item Consultation et comparaison des devis, arbitrages budgétaires.
  \item Réception des travaux ; réalisation moi-même d'une partie du chantier ; système vidéo et site web.
\end{itemize}

\entry{\textbf{ÉPOPÉE GADZ} --- trésorier, Europ'Raid 2026}{2025 -- 2026}
Budget et sponsoring d'un raid solidaire de 8\,000 km à travers l'Europe : près de 4\,000 euros récoltés.

\entry{\textbf{FunRide} --- technicien polyvalent, base de loisirs vélo \textnormal{; saisonnier aux vendanges}}{2020 -- 2021}
Réparation et entretien des vélos, accueil des clients, travaux manuels et logistique.

% ───── formation ─────
\section{Formation}

\entry{\textbf{Arts et Métiers (ENSAM)} --- \textit{Diplôme d'ingénieur Arts et Métiers (grade master)}}{2024 -- 2027}
Mécanique, résistance des matériaux, calcul par éléments finis, conception, procédés de fabrication.

\entry{\textbf{Classe préparatoire PT}, Lycée Mermoz, Montpellier --- \textit{maths, physique, sciences de l'ingénieur}}{2021 -- 2024}

% ───── projets ─────
\section{Projets}

\entry{\textbf{Pièce en fonte pour une fontaine d'Aix-en-Provence} --- équipe de 4 \quad\link{https://gauthierbonhomme.com/fonderie.html}{rapport}}{2025 -- 2026}
Simulation, moulage main, coulée ; imprimante en panne : adaptation à l'outillage existant.

\entry{\textbf{Calcul de structure sous Abaqus} --- équipe de 3 \quad\link{https://gauthierbonhomme.com/abaqus.html}{rapport}}{}
Modèle éléments finis d'une éprouvette fissurée : résultat à $\approx$\,1\,\% de la solution analytique.

\entry{\textbf{Vilebrequin --- de la gamme d'usinage à la pièce} \quad\link{https://gauthierbonhomme.com/usinage.html}{rapport}}{2025}
Gamme en FAO Fusion 360, production sur tour-fraiseur DMG MORI, chiffrage du coût.

\entry{\textbf{Atelier \& bricolage} --- personnel \quad\link{https://gauthierbonhomme.com/projet.html?id=projets-divers-et-varies}{photos}}{}
Moteurs remis à neuf, tables en chaudronnerie, jardinière, rénovation complète d'un vélo.

% ───── compétences ─────
\section{Compétences}

\textbf{Gestion de projet :} coordination des artisans, planning, comparaison de devis, réception de travaux, budget\\
\textbf{Technique :} résistance des matériaux, calcul éléments finis (Abaqus), lecture de plans, CATIA V5, 3DEXPERIENCE\\
\textbf{Outils :} Excel, Word, PowerPoint\\
\textbf{Savoir-être :} travail en équipe, rigueur, autonomie, adaptabilité\\
\textbf{Langues :} français (langue maternelle), anglais (C1)

% ───── sport ─────
\section{Sport \& montagne}

\textbf{Ski :} alpin et de randonnée\sep \textbf{Montagne :} alpinisme en autonomie, escalade (7B)\\
\textbf{Trail :} finisher d'un ultra-trail de 100 km (5\,000 m D+), prépare l'UTMB \quad\link{https://utmb.world/fr/runner/4895528.gauthier.bonhomme}{profil UTMB}\\
\textbf{Et aussi :} kitesurf\sep judo, ceinture noire

\end{document}
